\section{Methodology Detail, Per Sector}

This appendix is the detailed record of how each result was actually measured, sector by sector. For each result
it gives the exact method, the confound that had to be beaten, and the fix. The methodology chapter gives the
overview. This is the detail. Every result is tagged MEASURED for direct simulation, DERIVED for group theory or
exact calculation, or PREDICTED for a quantitative consequence not yet checked. The computational records are in
\cite{pollard2026vibetest}.

\subsection{Common methods}

These are used across sectors.

\begin{itemize}
  \item \textbf{Cell-graph construction.} The cell-graph builder builds the honeycomb's cell
  adjacency by tracking each cell as a group element, a matrix, and reaching neighbours via the face reflections.
  This is an $O(\mathrm{facet} \cdot n)$ breadth-first search. The flat-lattice builder
  builds the Euclidean references ($D_4$, $\mathbb{Z}^4$, $\mathbb{Z}^3$) directly. The hyperbolic builder stalls
  on flat honeycombs because their cell center is an ideal point, which is a correct signal of flatness.
  \item \textbf{BFS shells.} Breadth-first layers from a central cell give the growth profile, feeding dimension,
  curvature, cosmology, and hierarchy.
  \item \textbf{Lazy random-walk return probability.} $P(t) \sim t^{-d/2}$ gives the spectral dimension,
  $d = -2\, \mathrm{d}\log P / \mathrm{d}\log t$.
  \item \textbf{Bethe-lattice cavity recursion.} The exact infinite-tree Green's function, for gravity and
  holography.
  \item \textbf{Kernel Polynomial Method plus stochastic trace.} Spectral quantities such as sea energy and
  density of states, without diagonalization, matvec-only and GPU-scalable.
  \item \textbf{WebGPU.} Large-scale dynamics, the 24-direction gas, the bulk wave, the skyrmion, and the KPM, on
  the GPU.
\end{itemize}

\subsection{Gravity}

\begin{itemize}
  \item \textbf{Flat-3D Newtonian $1/r$, the difference method (p224).} The lattice Green's function of the
  Laplacian, $G(r)$. The naive finite-patch sum is contaminated by the $k = 0$ zero mode and boundary effects, and
  earlier gave a garbage exponent $-1.33$. The fix computes $G(r) - G(r_0)$ with the regular cosine integrand, so
  the $k = 0$ singularity cancels in the difference, then fits $G(r) - G(2) = a/r + b$. The result is
  $a \approx 0.083 \approx 1/(4\pi)$, a clean $1/r$. MEASURED.
  \item \textbf{Bulk-mediated correlator, the Bethe recursion (p239, p218).} The exact infinite-tree Green's
  function by the cavity recursion, with no finite patch. The boundary-to-boundary correlator decays as
  $1/r^\alpha$ with $\alpha = 2\ln(1/\mu)/\ln(b)$, and $\mu = 1/b$ universally, so $\alpha = 2$. Validated against
  a directly solved finite tree. DERIVED and MEASURED.
  \item \textbf{Cusp cross-check.} Re-measured on the $\{4,3,4\}$ cubic cusp by a screened-Laplacian solve.
  Spectral dimension is 3. The exponent there is small-$r$ and finite-size limited, about $0.67$, read as a
  convergence check. The clean $1/r$ is from the difference method.
  \item \textbf{Caveat.} The bulk $1/r^2$, the holographic 2-point function, and the flat-layer $1/r$, the
  Newtonian law, are different objects.
\end{itemize}

\subsection{Holography}

\begin{itemize}
  \item \textbf{Bulk-mediated boundary correlator (p239).} Bethe recursion, clean $1/r^2$. DERIVED and MEASURED.
  \item \textbf{Same correlator by tensor network (p241).} A tree-tensor (MERA) of random isometries, with the
  leaf-leaf 2-point function found by contracting the network. The decay is set by the ascending superoperator's
  second eigenvalue. The result is about $1/r^{1.93}$, matching Bethe by a different algorithm at polynomial cost,
  taming the exponential growth. Perfect tensors give the HaPPY code. MEASURED.
  \item \textbf{Hyperbolic band theory (p243).} On a closed hyperbolic manifold, the $\mathrm{PSL}(2,7)$ Cayley
  graph, the symmetry group is finite, so its irreps are the momenta and the spectrum block-diagonalizes by irrep.
  Eigenvalue degeneracies equal the irrep dimensions, $\{1, 6, 7, 8\}$, a momentum space with no real-space patch.
  MEASURED.
  \item \textbf{Emergent nonlocality, Bell (p236).} CHSH for a local deterministic strategy by brute force, max
  $2$, versus quantum with optimal angles, $2\sqrt{2} = 2.83$. A local discrete base obeys $\mathrm{CHSH} \le 2$,
  so quantum nonlocality requires a nonlocal boundary, which holography supplies, since two boundary points
  connect through the bulk. DERIVED.
  \item \textbf{Caveat.} Finite hyperbolic patches are about $99$ percent boundary, so naive bulk measurements are
  confounded. The holographic toolkit, recursion, closed manifolds, band theory, and tensor networks, is the
  correct method.
\end{itemize}

\subsection{Cosmology and expansion}

\begin{itemize}
  \item \textbf{Bulk shell growth (p237).} The cell-graph builder then a BFS sweep recording shell sizes, the
  cells at each radial distance, where radial distance equals cosmic time, one shell per beat. Shell sizes
  $1, 24, 456, 8376, 31143$ give a constant ratio $R \approx 11$ per beat, exponential. MEASURED.
  \item \textbf{Scale factor.} The 3D volume $\sim a^3$ grows by $R$ per beat, so $a(t) \sim R^{t/3} = e^{Ht}$ with
  $H = \ln(R)/3 \approx 0.80$ per beat, and $a''/a = H^2 > 0$, accelerating, with $\Lambda = 3H^2 \approx 1.9$ in
  beat units. DERIVED.
  \item \textbf{Flat control.} $\{3,4,3,3\}$ grows polynomially, ratio about $2.5$, with no de Sitter, so the
  expansion is a curvature effect.
  \item \textbf{Caveat.} The geometric ratio to physical Hubble rate needs a beat-to-second and cell-to-length
  identification, which is not pinned, so $H$ is in beat units. A multi-era cosmology with radiation and matter
  before de Sitter needs matter back-reaction, not yet modelled.
\end{itemize}

\subsection{Spin and fermions}

\begin{itemize}
  \item \textbf{Spinor coin (p190).} Decompose the 24-direction coin, the 24-cell equal to $D_4$, under the
  symmetry group, $24 = 8v + 8s + 8c$, the $\mathfrak{so}(8)$ vector plus two spinors, with $S_3$ triality
  permuting them, the three generations. The control $\{5,3,4\}$, 12 directions, gives $1 + 3 + 3' + 5$,
  icosahedral, no spinor. DERIVED.
  \item \textbf{Self $=$ fermion (p206).} A self is a hopfion, the Hopf map $S^3 \to S^2$, with charge in
  $\pi_3(S^2) = \mathbb{Z}$. Compute the Hopf linking charge, apply Finkelstein-Rubinstein and Wilczek-Zee, and an
  odd charge picks up $-1$ under $2\pi$, a fermion. Triple-confirmed by Hopf charge, configuration-space $\pi_1$,
  and spin-statistics. DERIVED.
  \item \textbf{Skyrme stabilizer (Skyrme-twist simulation, GPU).} Measure the fermion sea energy of a
  double-twist periodic texture by KPM, a Chebyshev expansion of $|H|$ with a stochastic trace, subtract a
  single-helix control by common random numbers to isolate the Skyrme term, and fit $A q^2 + B q^4$. The result is
  positive, $+798$ and $+814$, stabilizing, with exchange ratio about $1.5$. So matter forms and is stable, no
  fine-tuning. MEASURED.
  \item \textbf{Caveat.} Earlier localized-soliton attempts (p205, p214, p216) were confounded by
  surface-dominated sea energy. The double-twist periodic method with the helix control supersedes them.
\end{itemize}

\subsection{Gauge and the Standard Model}

\begin{itemize}
  \item \textbf{$\mathfrak{so}(10)$ from coin plus tone (p221).} Combine the $D_4$ root system,
  $\mathfrak{so}(8)$, with the ternary tone as an extra weight, building $D_5 = \mathfrak{so}(10)$. The 16 spinor
  is one generation, $16 = 8s$ (tone $+$) and $8c$ (tone $-$). DERIVED.
  \item \textbf{Breaking (p225, p227).} Take a self-condensate, the GUT Higgs, compute the unbroken subgroup,
  $\mathfrak{so}(10) \to \mathfrak{su}(5) \to \mathrm{SM}$, the Georgi-Glashow chain, with
  $16 = 1 + 10 + \bar{5}$. Any self-condensate gives $\mathfrak{su}(5)$, 20 unbroken roots, so the breaking is
  automatic. DERIVED.
  \item \textbf{$\sin^2\theta_W = 3/8$ (p228).} Sum over one generation,
  $\sin^2\theta_W = \sum T_3^2 / \sum Q^2 = 3/8$ exactly at unification. DERIVED. The running to experiment is in
  the elevating-the-predictions appendix.
  \item \textbf{Mass relations (p231).} $b$-$\tau$ unification, $m_b = m_\tau$ at the GUT scale, the determinant
  relation $\det(M_{\mathrm{lepton}}) = \det(M_{\mathrm{down}})$ from the discrete $\mathrm{Tr}\,Y = 0$ over the
  16, and the Georgi-Jarlskog factor-3 relations. DERIVED.
  \item \textbf{Non-abelian gauge field (p234).} $SO(3)$ link matrices, the prototype, identical in form for
  $\mathfrak{so}(10)$. The Wilson loop, the trace of the ordered product, is gauge-invariant, has non-trivial
  holonomy (curvature), and the links do not commute (non-abelian). MEASURED. The symmetric collision making
  $\mathfrak{so}(10)$ the local symmetry is IR-forced by triality (p226).
\end{itemize}

\subsection{Dirac and Lorentz}

\begin{itemize}
  \item \textbf{Dirac dispersion (p230).} Simulate the discrete two-component left and right mover walk, measure
  $E(k)$ by Fourier-transforming a tracked cell's time series. Massless, the pure discrete shift gives $E(k) = k$
  exactly, no continuity. Massive, add a coin rotation by mass angle $m$, get $\cos E = \cos(m)\cos(k)$, the
  lattice Dirac, with rest mass $E(0) = m$, so mass is the coarse-grained mixing rate. MEASURED.
  \item \textbf{Light cone (p193, p211).} Propagate a pulse, front speed one cell per beat ($z = 1$), a finite
  invariant speed. MEASURED.
  \item \textbf{Caveat.} p230 is the clean 1D two-component demonstrator. The full 24-direction version is the bulk
  wave, and cusp Lorentz invariance carries the small cubic anisotropy below, isotropic in the IR. The complex
  quantum amplitudes are the emergent description, the base is the real discrete walk.
\end{itemize}

\subsection{Isotropy}

\begin{itemize}
  \item \textbf{Dispersion isotropy (p233, analytic).} Compute
  $\omega^2(k) = \sum_{\mathrm{dir}} (1 - \cos(k \cdot d))$ for the $D_4$ 24-direction coin versus cubic
  $\mathbb{Z}^3$ and hypercubic $\mathbb{Z}^4$, compare axis versus body diagonal at the same $|k|$, and check the
  fourth-moment condition $\sum d_i^4 = 3 \sum d_i^2 d_j^2$. $D_4$ has an isotropic fourth moment, $12 = 3 \cdot
  4$, isotropic to order 4, zero anisotropy, versus about 8 to 9 percent for cubic. MEASURED.
  \item \textbf{Light-cone isotropy (isotropy simulation, GPU).} Pulse on a 4D lattice with the 24 $D_4$
  directions versus the 8-direction hypercubic, measure front shape. The 24-direction front is round, ratio
  $1.00$, versus faceted, $0.50$, on a 4D lattice of 331k cells. MEASURED.
  \item \textbf{Symmetry restoration (p226).} Invariant-polynomial dimensions for $B_4$ (hypercubic) versus $F_4$
  (24-cell with triality), via Newton's identities on power traces. $F_4$ has one degree-4 invariant versus
  $B_4$'s two, so triality kills the order-4 anisotropy, discrete $F_4 \to$ continuous $\mathfrak{so}(8)$ in the
  IR. DERIVED.
  \item \textbf{Caveat.} This is the bulk $D_4$ isotropy. The physical cusp is the $\{4,3,4\}$ cubic, anisotropic
  at order 4, the small $(E/E_{\mathrm{Planck}})^2$ residual, so matter-propagation isotropy on the cusp is the
  cubic one, IR-isotropic with a tiny UV anisotropy. This bulk-versus-cusp distinction is the key subtlety.
\end{itemize}

\subsection{Electromagnetism}

\begin{itemize}
  \item \textbf{U(1) Gauss law (p223).} The rule's local charge conservation, charge in equals charge out per cell
  per beat, is a discrete Gauss law, the conserved current sourcing the emergent U(1). A direct conservation check
  on streaming. MEASURED.
  \item \textbf{Emergent U(1) gauge field (p232).} A vortex gauge field, link phases winding around a plaquette to
  give flux $\Phi$, on a 2D lattice. First, the Wilson loop, the sum of link phases, equals the enclosed flux and
  is gauge-invariant under a random $A \to A + \mathrm{d}\lambda$, flux $0.7$ unchanged. Second, Aharonov-Bohm, a
  charge encircling the flux picks up exactly $\Phi$, loop- and gauge-independent, and a non-enclosing loop gives
  $0$. Third, the static sector is Coulomb $1/r$, the flat-3D Green's function (p224). MEASURED.
  \item \textbf{Caveat.} A single link gives a flux dipole, net zero. The fix was the vortex winding construction
  giving a single-plaquette flux, which made the Wilson loop clean. The non-abelian extension is the same
  construction with matrices (p234).
\end{itemize}

\subsection{Physical space and the cusp}

\begin{itemize}
  \item \textbf{Horosphere flatness (horosphere-reality check).} Extract a horosphere band from a finite
  $\{3,4,3,4\}$ ball, a Busemann-pruned BFS slab from the horosphere-band builder, measure intrinsic growth
  (polynomial means flat) and degree histogram (uniform means clean periodic). The result is intrinsically flat at
  finite distance, growth $1.52$, not the bulk's roughly $9$. MEASURED.
  \item \textbf{Clean cusp versus generic slice (emergent-space check).} On the $\{3,4,3,4\}$ horosphere
  band and the clean $\{4,3,4\}$ cubic cusp, measure spectral dimension, light-cone isotropy (front-radius CV),
  and connectivity. The clean cusp is genuine 3D, spectral dimension $2.97$, while a generic slice is rough and
  sub-3D, $2.16$. MEASURED.
  \item \textbf{Sufficient size (cusp-convergence check).} The cusp's intrinsic structure is the $\{4,3,4\}$
  cubic lattice, cheap $O(L^3)$ $\mathbb{Z}^3$ boxes. Build boxes of growing side $L$ and watch spectral dimension
  approach 3 and the gravity exponent become $L$-independent. Continuum-like by $L \approx 17$ cells, a few beats,
  and the observed universe at about $10^{61}$ cells is hugely sufficient. MEASURED.
  \item \textbf{Physics holds on real space (physics-on-real-space check).} Compute the flat-layer observables
  on the $\{4,3,4\}$ cusp versus a generic aperiodic slice. The physics holds on the cusp cubic and degrades on a
  generic slice, so physical space equals the cusp cubic. MEASURED.
  \item \textbf{Bulk-to-cusp projection (bulk-cusp-interface check).} Verify the spinor branching
  $SO(4) \to SO(3)$, 4D Dirac to 3D Pauli, $4 \to 2 + 2$, and project the 24 bulk directions onto the cusp
  tangent. DERIVED and MEASURED.
  \item \textbf{Caveat.} A generic Busemann slice is a bad extraction, thin, aperiodic, and sub-3D. This corrected
  an earlier optimistic guess that aperiodic meant isotropic. Whether the growing mesh dynamically settles matter
  into clean cusp regions is the one open dynamical question, a defined next GPU experiment.
\end{itemize}

\subsection{Substrate selection}

\begin{itemize}
  \item \textbf{Sweep every dimension.} Build each regular honeycomb with the cell-graph builder, plus
  the flat-lattice builder for flat references, measure coin degree, spectral dimension, the
  crystallographic test (only 3, 4, 6 in the symbol), growth ratio (curvature), compactness class, and the Bethe
  holographic exponent. Coverage is 8 in 2D, 15 in 3D (4 compact plus 11 paracompact), the full 4D set, and the 5D
  candidates plus $\{5,3,3,3,3\}$.
  \item \textbf{The dimension ladder.} Physical-space dimension equals bulk dimension minus one, since the
  horosphere is one dimension lower. 2D to 1D, 3D to 2D, 4D to 3D, 5D to 4D. Only 4D gives the observed 3D.
  MEASURED.
  \item \textbf{Same-battery comparison (cross-dimension comparison).} The identical full battery on one
  representative per dimension, $\{7,3\}$, $\{5,3,4\}$, $\{3,4,3,4\}$, $\{5,3,3,3,3\}$, plus the flat twin
  $\{3,4,3,3\}$, apples-to-apples.
  \item \textbf{Scored ranking (tessellation ranking).} Score by a rubric, spinor coin 3,
  crystallographic 2, physical-dimension closeness 3, curvature 1, compactness, and $\{3,4,3,4\}$ tops it at
  $9.5$. DERIVED and MEASURED.
  \item \textbf{Caveats.} High-degree and 120-cell honeycombs saturate a finite patch, giving artifact
  spectral-dimension and growth, so the clean reads are the low-degree ones. The lower-D substrates are not
  strictly ruled out, since their missing features could emerge at higher layers. $\{3,4,3,4\}$ is the
  least-emergence choice, not the only conceivable one. The rubric weights are a judgment call, but $\{3,4,3,4\}$
  tops any reasonable weighting because it alone is non-zero on every factor.
\end{itemize}
